\section{Conditional path laws at the same roof value}
\label{sec:v52-same-roof}
We now pass from a path-valued raw numerator to a conditional path
kernel at its own roof value. Convergence is first proved in mean
with respect to the unchanged conditional roof law. A second result
shows that the missing scalar positive-height estimate would also
supply the uniform pointwise path bridge, without a separate
uncontrolled bounded-Lipschitz numerator.

Fix $h,d>0$ and $M_0<\infty$. Put $I_t=[t,t+h]$ and retain the
central coordinate $Z(t)$ of Section~\ref{sec:v51-positive-remainder}.
Our uniform target class is
\begin{equation}\label{eq:v52-target-class}
 |Z(t)|\le M_0,\qquad
 H_m:=\int_{I_t}(G_{m,R}^{n,k}(u))_+\,du\ge dh.
\end{equation}
Only an integrated lower bound is required here. Let
$F_m=\int_{I_t}P_{m,R}^{n,k}(u)\,du$. The scalar local-variation
law gives $F_m\ge dh/2$ eventually. The original roof law under
$E=\{K_{n,R}=k,N_{n,R}=m,T_{n,R}\in I_t\}$ and its reference are
\[
 \pi_m(du)=\frac{\one_{I_t}(u)P_{m,R}^{n,k}(u)}{F_m}\,du,
 \qquad
 \widehat\pi_m(du)=\frac{\one_{I_t}(u)(G_{m,R}^{n,k}(u))_+}{H_m}\,du.
\]
Their probability total-variation distance is one half the variation
mass and tends uniformly to zero by
Corollary~\ref{cor:v49-conditional-roof-TV}.

\subsection{Return-clock transfer in the roof-integrated path norm}
Let $\mathsf R_{m,R}^{n,k,u}$ denote the conditional law of the actual
return bridge $\mathcal Y_{n,R}$ under the same source disintegration
used to define $\mathsf Q_{m,R}^{n,k,u}$. Its reference is
$\mathsf V_R=\operatorname{Law}(\mathbb B_{D_R})$.

\begin{lemma}[Integrated transfer of the pinned clock]
\label{lem:v52-return-transfer}
For every fixed $h,M_0$, uniformly over $|Z(t)|\le M_0$, without
requiring a positive denominator,
\begin{equation}\label{eq:v52-return-path-variation}
 \frac1h\int_{I_t}
 \left\|P_{m,R}^{n,k}(u)\mathsf R_{m,R}^{n,k,u}
                  -G_{m,R}^{n,k}(u)\mathsf V_R\right\|_{\mathrm{BL}^*}
                                                       \,du\longrightarrow0.
\end{equation}
\end{lemma}
\begin{proof}
On $E$ one has $n/m=c+O_{M_0}(m^{-1/2})$. Put
$\mathsf A_{m,R}=\sqrt{m/n}\,L_R^{-1}$, where $L_R$ is the covariance change
matrix in \eqref{eq:v37-covariance-change}. These matrices are uniformly
bounded, and
$(\mathsf A_{m,R})_*\mathsf W_R\to\mathsf V_R$ uniformly in bounded-Lipschitz
distance, by covariance continuity and $m/n\to c^{-1}$.

Consider the finite, unnormalized source measure
$\lambda_m=(m^2/h)\one_E\nu_R^*$. Its mass is bounded uniformly by
the scalar fixed-window upper bound. The unnormalized fourth moment
\eqref{eq:v41-pinned-fourth} gives a uniformly tight family of
collision-path measures under $\lambda_m$, just as in
\eqref{eq:v52-band-path-fourth}. This argument is valid when the
mass tends to zero: add mass at the zero path and normalize by a
common upper bound. No rare-event denominator is introduced.

At return $l$ put $\theta_l=N_{l,R}/m$. The exact half-open
occupation convention gives $A_{N_{l,R},R}=l$ and $A_{m,R}=n$.
The fourth collision-bridge coordinate consequently satisfies
\begin{equation}\label{eq:v52-exact-clock-displacement}
 \theta_l-l/n=-\frac{\sqrt m}{n}
                  (\mathcal B_{m,R})_4(\theta_l).
\end{equation}
Tightness makes the maximum of these displacements tend to zero in
$\lambda_m$-measure. The centered terminal vector
$\mathcal X_{m,R}(1)$ is bounded on $E$, because the actual roof lies
in $I_t$ and $|Z(t)|\le M_0$; explicitly,
$|Z(u)|\le M_0+h/\sqrt m$ on that interval. The exact pinned time-change identity
\eqref{eq:v41-pinned-time-change}, followed by the uniform modulus
of continuity of a tight family of continuous paths, now gives
\begin{equation}\label{eq:v52-clock-coupling}
 \int \min\{2,\|\mathcal Y_{n,R}-\mathsf A_{m,R}\mathcal B_{m,R}\|_\infty\}
                                                     \,d\lambda_m\to0.
\end{equation}
Linear interpolation on the return grid is paid by the same modulus.
Neither a changed roof event nor a change of return index occurs.

Use the joint source disintegration to couple the two conditional
paths at each individual roof. Integrating their bounded-Lipschitz
distance, with weight $P/h$, is bounded by
\eqref{eq:v52-clock-coupling}. Under a linear map $A$, the path dual
norm grows by at most $\max\{1,\|A\|\}$. Apply this to
\eqref{eq:v52-path-local-variation}. Finally $G$ is uniformly bounded
and the reference covariance convergence pays
$G[(\mathsf A_{m,R})_*\mathsf W_R-\mathsf V_R]$ in the integrated norm.
These three estimates prove \eqref{eq:v52-return-path-variation}.
\end{proof}

\begin{theorem}[A same-roof bridge in conditional mean]
\label{thm:v52-same-roof-mean}
Uniformly on the class \eqref{eq:v52-target-class},
\begin{align}
 \int_{I_t}d_{\rm BL}(\mathsf Q_{m,R}^{n,k,u},\mathsf W_R)
                                                   \,\pi_m(du)&\to0,
                                             \label{eq:v52-collision-kernel-mean}\\
 \int_{I_t}d_{\rm BL}(\mathsf R_{m,R}^{n,k,u},\mathsf V_R)
                                                   \,\pi_m(du)&\to0.
                                             \label{eq:v52-return-kernel-mean}
\end{align}
Both conclusions also hold with $\widehat\pi_m$ in place of $\pi_m$.
The conditional kernels fix the original exact $n,k,m$ and their own
roof value $u$. They are not the path law under a substituted window.
No polynomial rate in $m$ is asserted.
\end{theorem}
\begin{proof}
For a nonnegative scalar $P$, a real scalar $G$, and probabilities
$Q,W$, the constant test one belongs to $\mathcal F$, so
\[
 P\|Q-W\|_{\mathrm{BL}^*}
 \le\|PQ-GW\|_{\mathrm{BL}^*}+|P-G|
 \le2\|PQ-GW\|_{\mathrm{BL}^*}.
\]
Integrate this inequality over $I_t$ and divide by $F_m\ge dh/2$.
Theorem~\ref{thm:v52-path-remainder} proves the collision assertion;
Lemma~\ref{lem:v52-return-transfer} proves the return assertion.
The integrands are measurable and bounded by two. Replacing the
roof measure by $\widehat\pi_m$ changes their integrals by at most
four times the probability total-variation distance, which tends
to zero. This uses the actual path numerators, not merely the
scalar roof approximation.
\end{proof}

\begin{corollary}[Asymptotically full sets of roof values]
\label{cor:v52-full-roof-sets}
There is a deterministic sequence $a_m\downarrow0$, depending on
$h,d,M_0$, such that outside target-dependent Borel sets of both
$\pi_m$- and $\widehat\pi_m$-mass tending uniformly to zero, the sum
of the two distances in
\eqref{eq:v52-collision-kernel-mean}--\eqref{eq:v52-return-kernel-mean}
is at most $a_m$. The sequence is qualitative, not a prescribed
shrinking roof resolution. No assertion is made at every externally
chosen roof value.
\end{corollary}
\begin{proof}
Let $\gamma_m$ be the supremum over targets of the two expectations
under both roof laws. It tends to zero. Replace it by its decreasing
tail supremum and take, for example,
$a_m=\sqrt{\sup_{j\ge m}\gamma_j}+m^{-1}$.
Markov's inequality bounds each exceptional mass by
$\gamma_m/a_m\to0$. Kernel versions can first be fixed on the
standard Borel spaces; all source-null modifications are immaterial.
\end{proof}

\subsection{The scalar height condition also controls the path numerator}
On the set $\{G_{m,R}^{n,k}\ge d\}$ the lower law ensures
$P\ge d/2$ eventually. There the exact physical-defect posterior is
\begin{equation}\label{eq:v52-physical-posterior}
 \beta_{B,m,R}^{n,k}(u)=\frac{\mathfrak b_{B,m,R}^{n,k}(u)}
                              {P_{m,R}^{n,k}(u)}
 =\mathbb E_{\nu_R^*}[1-H_{m,R}^{\varepsilon(B)}
                       \mid \Pi_{n,k,m,R},T_{n,R}=u]\in[0,1].
\end{equation}
This is a regular conditional expectation, not a positive probability
for the equality $T_{n,R}=u$.

\begin{proposition}[Pointwise path error charged to the same physical source]
\label{prop:v52-pointwise-path-charge}
For every fixed sufficiently large $B$,
\begin{equation}\label{eq:v52-pointwise-path-charge}
 \limsup_{m\to\infty}\sup_{R,n,k}
 \operatorname*{ess\,sup}_{u:G(u)\ge d}
 \bigl(d_{\rm BL}(\mathsf Q_{m,R}^{n,k,u},\mathsf W_R)
                       -2\beta_{B,m,R}^{n,k}(u)\bigr)_+
                                              \le C_dB^{-1/192}.
\end{equation}
On a fixed central compact set, the scalar positive-height condition
\eqref{eq:v51-positive-criterion} therefore implies
\begin{equation}\label{eq:v52-pointwise-bridge-implication}
 \sup_{R,n,k}\operatorname*{ess\,sup}_{\substack{u:\ |Z(u)|\le M_0\,,\ G(u)\ge d}}
 d_{\rm BL}(\mathsf R_{m,R}^{n,k,u},\mathsf V_R)\longrightarrow0.
\end{equation}
It also implies the corresponding collision-bridge statement. This
is an implication from the still-unproved height condition, not an
unconditional uniform pointwise bridge theorem.
\end{proposition}
\begin{proof}
Put $\mathsf E=\mathbf P-G\mathsf W_R-\mathbf b_B$.
Taking masses gives $P-G-\mathfrak b_B=\mathsf E(1)$.
Consequently
\[
 P(\mathsf Q-\mathsf W_R)
     =\mathbf b_B-\mathfrak b_B\mathsf W_R
                         +\mathsf E-\mathsf E(1)\mathsf W_R.
\]
Its dual norm is at most $2\mathfrak b_B+2\|\mathsf E\|_{\mathrm{BL}^*}$.
Divide by $P\ge d/2$ and apply
Theorem~\ref{thm:v52-path-remainder}. This proves
\eqref{eq:v52-pointwise-path-charge}.
If \eqref{eq:v51-positive-criterion} holds, the same division shows
that the conditional collision-bridge error tends uniformly to zero
on the displayed positive-reference central sets, taking the collision
limit before the reconstruction band tends to infinity.

For the return conclusion take any hypothetical violating sequence
of such targets and nonexceptional roof values. After a subsequence
$R\to R_0$. The just-proved conditional collision-path convergence
gives tightness for this sequence of conditional laws. The exact
identity \eqref{eq:v52-exact-clock-displacement} and the bounded
central terminal vector give the same pinned clock comparison as in
Lemma~\ref{lem:v52-return-transfer}, now in those individual
conditional laws. The deterministic linear covariance change tends
to $D_{R_0}$. Thus the conditional return laws converge to
$\mathsf V_{R_0}$, contradicting the sequence. This proves
\eqref{eq:v52-pointwise-bridge-implication} in essential supremum.
The new numerator estimate removes the need for a separate bounded-
Lipschitz path remainder condition; it does not prove the remaining
scalar incidence or clearance height bounds.
\end{proof}

\paragraph{The remaining pointwise endpoint.}
Theorems~\ref{thm:v52-path-remainder} and
\ref{thm:v52-same-roof-mean} give a path-valued local density law and
same-roof conditional convergence in roof probability on the full
original source. Their path topology is weak, expressed by
bounded-Lipschitz distance; it is not path-space total variation.
The scalar two-sided essential-height theorem still requires
\eqref{eq:v51-positive-criterion}. Its positive incidence and clearance
components have not been shown small in height here. The arithmetic
transition kernel remains in every uniform assertion; no residue is
set equal to one. The fixed-radius positive-class specialization uses
precisely the unchanged arithmetic lower bounds already established
in Section~\ref{sec:v51-roof-minorization}.
